% !TeX program = pdflatex
\documentclass[12pt,a4paper]{article}

\newcommand{\notelanguage}{finnish}
% Shared preamble for course notes and exercise sets.

\usepackage[T1]{fontenc}
\usepackage{lmodern}
\usepackage[english]{babel}

\usepackage[a4paper,margin=30mm]{geometry}
\usepackage{microtype}
\usepackage{graphicx}
\usepackage{enumitem}

\usepackage{amsmath}
\usepackage{amssymb}
\usepackage{amsthm}
\usepackage{mathtools}
\usepackage{aliascnt}

\usepackage[hidelinks,pdfusetitle]{hyperref}

\numberwithin{equation}{section}
\setlist{itemsep=0.25em,topsep=0.5em}

\theoremstyle{plain}
\newtheorem{theorem}{Theorem}[section]
\newaliascnt{lemma}{theorem}
\newtheorem{lemma}[lemma]{Lemma}
\aliascntresetthe{lemma}
\newaliascnt{proposition}{theorem}
\newtheorem{proposition}[proposition]{Proposition}
\aliascntresetthe{proposition}
\newaliascnt{corollary}{theorem}
\newtheorem{corollary}[corollary]{Corollary}
\aliascntresetthe{corollary}

\theoremstyle{definition}
\newaliascnt{definition}{theorem}
\newtheorem{definition}[definition]{Definition}
\aliascntresetthe{definition}
\newaliascnt{example}{theorem}
\newtheorem{example}[example]{Example}
\aliascntresetthe{example}
\newaliascnt{problem}{theorem}
\newtheorem{problem}[problem]{Problem}
\aliascntresetthe{problem}

\theoremstyle{remark}
\newaliascnt{remark}{theorem}
\newtheorem{remark}[remark]{Remark}
\aliascntresetthe{remark}

\usepackage[nameinlink,noabbrev]{cleveref}

\crefname{theorem}{theorem}{theorems}
\crefname{lemma}{lemma}{lemmas}
\crefname{proposition}{proposition}{propositions}
\crefname{corollary}{corollary}{corollaries}
\crefname{definition}{definition}{definitions}
\crefname{example}{example}{examples}
\crefname{problem}{problem}{problems}
\crefname{remark}{remark}{remarks}

\DeclareMathOperator{\im}{im}
\DeclareMathOperator{\rank}{rank}
\DeclarePairedDelimiter{\abs}{\lvert}{\rvert}
\DeclarePairedDelimiter{\norm}{\lVert}{\rVert}
\DeclarePairedDelimiter{\set}{\lbrace}{\rbrace}


\usepackage{tasks}

\settasks{
  label = (\alph*),
  label-width = 2em,
  after-item-skip = 0pt
}

\title{Harjoituksen 2 ratkaisut}
\date{}
\author{}

\begin{document}

\maketitle

\noindent

\begin{exercise}
  Kirjoita käyttäen matemaattista notaatiota (joukko-opillisia ja loogisia
  merkkejä)
\begin{enumerate}[label=\alph*), nosep]
  \item luku \(\sqrt{2}\) ei kuulu rationaalilukuihin,
  \item lukujen \(-7\), \(0\), \(51\) ja miljoona muodostama joukko on
    kokonaislukujen osajoukko,
  \item joukko \(C\) on joukkojen \(A\) ja \(B\) leikkauksen komplementin
    aito osajoukko,
  \item niiden reaalilukujen \(x\) joukko, joiden neliö on suurempi kuin
    \(1{,}2\) ja pienempi kuin \(3{,}6\),
  \item lukujen \(a\) ja \(b\) irrationaalisuudesta ei seuraa summan
    \(a + b\) irrationaalisuus,
  \item joukkojen \(A\) ja \(B\) yhdiste sisältää joukkojen \(C\) ja \(D\)
    leukkauksen,
  \item jokaista nollasta poikkeavaa kokonaislukua \(x\) kohti on olemassa
    sellainen luku \(y\), että \(xy = 2\),
  \item \(\RR\):n suljetun välin kolmesta viiteen ja avoimen välin neljästä
    kuuteen joukkoerotus on suljettu väli kolmesta neljään.
\end{enumerate}
\end{exercise}

\begin{solution}
\leavevmode\par
\begin{enumerate}[label=\alph*), nosep]
\item \(\sqrt{2} \notin \QQ\)
\item \(\set{-7, 0, 51, 10^6} \subseteq \ZZ\)
\item \(C \subset \complement{(A \cap B)}\)
\item \(\set{x \in \RR \colon 1{,}2 < x^2 < 3{,}6}\)
\item \(\bigl((a \in \RR \setminus \QQ) \land (b \in \RR \setminus \QQ)\bigr)
  \nRightarrow a + b \in \RR \setminus \QQ\)
\item \(C \cap D \subseteq A \cup B\)
\item \(\forall x \in \ZZ \setminus \set{0} \; \exists y \in \RR
  \colon xy = 2\)
\item \([3, 5] \setminus \left]4, 6\right[ = [3, 4]\)
\end{enumerate}
\end{solution}

\begin{exercise}
Olkoon \(X = \set{a, b, c, d, e}\) perusjoukko sekä
\(A = \set{a, b, c}\) ja \(B = \set{b, c, d}\).
Määritä
\begin{enumerate}[label=\alph*), nosep]
\item \(A \cup B\) ja \(A \cap B\),
\item \(\overline{A}\) ja \(\overline{B}\),
\item \(A\):n ja \(B\):n \emph{symmetrinen erotus} \(A \Delta B = (A
\setminus B) \cup (B \setminus A)\),
\item \(B\):n \emph{potenssijoukko} \(\powerset{B} = \set{C : C
\subseteq B}\).
\end{enumerate}
Havainnollista joukkoja Vennin diagrammeilla.
\end{exercise}

\begin{solution}
\end{solution}

\begin{exercise}
Olkoon \(A\) ja \(B\) joukkoja. Todista \emph{absorbiolait}
\begin{enumerate}[label=\alph*), nosep]
\item \(A \cup (A \cap B) = A\),
\item \(A \cap (A \cup B) = A\).
\end{enumerate}
Havainnollista tuloksia Vennin diagrameilla.
\end{exercise}

\begin{solution}
\end{solution}

\begin{exercise}
Olkoon \(A_k = \left]0, \frac{1}{k}\right[, k \in \NN \setminus \set{0}\).
Määritä joukot
\begin{tasks}(4)
\task \(\displaystyle \bigcup_{k=1}^n A_k\),
\task \(\displaystyle \bigcap_{k=1}^n A_k\),
\task \(\displaystyle \bigcup_{k=1}^{\infty} A_k\),
\task \(\displaystyle \bigcap_{k=1}^{\infty} A_k\).
\end{tasks}
\end{exercise}

\begin{solution}
\end{solution}

\begin{exercise}
Kirjoita määritelmän avulla inkluusio
\[
A \cup (B \cap C) \subseteq (A \cap B) \cup (A \cap C)
\]
implikaationa ja todista se vastatodistuksen avulla.
\end{exercise}

\begin{solution}
\end{solution}

\begin{exercise}
Olkoon \(A\), \(B\) ja \(C\) joukkoja. Todista karteesisen tulon osittelulaki
leikkauksen suhteen:
\[
(A \cap B) \times C = (A \times C) \cap (B \times C).
\]
Havainnollista tulosta joukoilla \(A = \set{1, 2, 3}\), \(B = \set{2, 3, 4}\)
ja \(C = \set{1, 2}\) määrittämällä ensin joukko \(A \cap (B \times C)\)
ja sitten joukko \((A \times C) \cap (B \times C)\).
\end{exercise}

\begin{solution}
\end{solution}

\end{document}
